%!TEX root =  ../article.tex
%%%%%%%%%%%%%%%%%%%%%%%%%%%%%%%%%%%%%%%%%%%%%%%%%%%%%%%%%%%%%%%%%%%%%%%%%%%%%%
\section{Introduction}
%%%%%%%%%%%%%%%%%%%%%%%%%%%%%%%%%%%%%%%%%%%%%%%%%%%%%%%%%%%%%%%%%%%%%%%%%%%%%%

% CONTEXT

\emph{Determinacy races} are concurrent programming hazards occurring when two accesses on the same memory address are not ordered, and at least one is writing~\cite{determinator}.
The presence of determinacy races in a program makes its output varying given the same input, which may be hint the presence of correctness errors.
With ever growing intra-node parallelism of supercomputers, porting simulation codes while ensuring parallel execution correctness is challenging.
It is particularly true with asynchronous and dependent task-based programming models, where a missing synchronization can lead to an incorrect order of execution.

MOTIVATIONS

\romain{Motivations here !! ROMP also uses DBI ... so gotta motivates differently !! probably around false-negative}

Existing determinacy race detectors~\cite{archer,Matar18} relies on LLVM' ThreadSanitizer~\cite{thread-sanitizer-llvm}.

In the specific case of Archer~\cite{archer}, 

Motivations and results of Archer~\cite{archer} showed the effectiveness (even at large scale) of such approach for debugging applications.
However, as suggested by E. Dijkstra~\cite{EWD303}, \emph{"debugging is an inadequate means for ensuring programs correctness and that we must prove the correctness of programs"}.
 

It requires the entire swith instrumentation enabled, otherwise, analysis may miss races occuring in uninstrumented code.

OBJECTIVES

This paper explores the use of heavyweight dynamic binary instrumentation towards no false-negative determinacy race analysis. Our contribution are:
\begin{itemize}
	\item	the introduction of Taskgrind: a Valgrind tool for parallel programming model memory accesses analysis.
	\item 	a survey on pitfalls related to heavyweight DBI in parallel programs,
	\item 	the evaluation of a determinacy race analysis implemented into Taskgrind, comparing it to state of the art tools.
\end{itemize}

The paper is organized as follows.
\sect{sec:background} provides background on races analysis and Valgrind.
\sect{sec:taskgrind} introduces the design of Taskgrind and the determinacy race detector implementation.
\sect{sec:pitfalls} presents a few pitfalls on parallel programs heavyweight DBI, and how Taskgrind workarounds them.
\sect{sec:eval} evaluates Taskgrind overheads and its determinacy race detection analysis.
Finally, \sect{sec:relworks} review the literature and we conclude in \sect{sec:conclusion}.